\section{章节罗盘：危机怎样进入财政、家庭与口岸}
\label{sec:mid-qing-compass}

本章从盛世扩张的后果出发，而不预设一个突然坍塌的帝国。战争善后、旗人生计、河工漕运、盐课、灾荒与边务同时争取有限的钱粮和行政注意；它们分别落在京师账簿、地方道路和家庭借贷之中。危机因而不是一项可由总数概括的衰败，而是多种制度承诺开始彼此挤压。

白莲教后的恢复、边地行动与沿海贸易会显示这种挤压如何累积。鸦片、禁烟和战争随后将旧有的财政与社会难题接入条约、赔款和口岸程序。本章并不把对外冲突当作唯一原因；它要追问的是，清廷为何仍能采取行动，却越来越难让一项处置同时在前线、州县和民户那里奏效。这一失配将成为十九世纪战争国家的起点。
